\quellseite{72}
\part[Part III. The Catechist’s Text-Book in History]{Part III. Historical and Critical—The Catechist’s Text-Book in History}

\quellseite{74}
\chapter[Chapter VI. Catechisms Before Luther]{Chapter VI. The Catechist’s Text-Book. A Brief History of Catechisms Before Luther}

In the former chapters we have looked into the history of catechising in the pre-Reformation Church. We saw that it had its lights and its shades, its periods of bloom and promise as well as its periods of decay. The variations depended mainly on the Catechists. Where we find a teacher with the proper qualifications there we find good catechising. Where he is defective or positively bad there the work suffers and decays.

In this chapter we wish to review the helps and text-books that the Catechist has used in the past ages of the Church.

The task of the teacher is to make known and apply the truth. What truth? The truth that the soul needs. The truth that concerns man’s higher nature and eternal destiny. The truth that satisfies the deepest longings of his innermost spirit. \randnotiz{The truth to be taught.} God’s truth. The truth which He alone has in Himself and which He alone can impart, which endureth forever, is a shield and buckler, lived on by Old Testament saints, personified in Him who said, “I am the truth,” whose Spirit is the Spirit of truth
\quellseite{75}
which sanctifies, makes wise unto salvation, whose knowledge is life eternal.\randnotiz{The truth in the Word.} That truth is found only in the Word of God and in teachers and books that have drawn it out of that living and life-giving Word.

The teacher, then, is to impart the truth which God has revealed in His Word. But how? In what order? By what system? The Bible itself is not given in the form of a text-book. Its teachings are not put into the form of a finished pedagogical system. God has given the truth of Revelation as He has given the truth of physical nature. The truths of astronomy, of geology, of botany, or of zoölogy are all in nature. But it would be hard for an ignorant beginner to learn these sciences with no teacher or book except nature. God has given to man a wonderful mind. With this mind God expects man to dig out and arrange truth for his own enlightenment and use. And so out of nature man has formulated the natural sciences and written his text-books.


And so God expects man to dig out, arrange and systematize the truth that is contained in His Word. \randnotiz{This truth to be arranged and adapted.} The teacher must have a system, a plan and helps for his pupils. He must know how to arrange and adapt his teaching to his pupils. He must give milk to the babes and strong meat to the mature.


To thus adapt the truth he must understand his

\quellseite{76}
pupils, their whole makeup, physical, mental and spiritual. He must know something of the growth and development of the powers and capacities of his pupils. Knowing their psychology he must know enough of the science and art of pedagogy to adapt the truth as the pupil can apprehend it. To this we shall return later.

The early Church teachers did not have either the psychology or the pedagogy developed as we have it now. They had neither a Lutheran Graded Sunday School Series, nor a Luther’s Small Catechism. But the earnest teachers were feeling after a development with such help as they could. \randnotiz{Early Church and psychology and pedagogy.}

We have seen above that Justin Martyr shows how the candidates for baptism were instructed in the rudiments of evangelical teaching. So in the Catechetical Sermons of Cyril of Jerusalem, and in the Apostolic Constitutions we have seen the range of the teaching enlarged. In Irenæus (died 202) we find an incomplete form of the Apostles’ Creed for catechumens. This was orally explained and discussed. Tertullian wrote an exposition of the Lord’s Prayer, in which he also emphasizes the duty of prayer in general (“De Oratione”). So also Cyprian. Rufinus wrote an explanation of the Apostles’ Creed (“Expositio Symbolum Apostolicum”). Likewise Augustine (“De Fide et Symbolo”), with frequent reference to the Commandments.

\quellseite{77}
These and other early writings show us how the Fathers attempted a systematizing of revealed truth. They show further that their catechumens were instructed, even as ours are, in the Ten Commandments, the Lord’s Prayer and the Creed. Of necessity they would be instructed also concerning baptism, before they received it, and concerning the Lord’s Supper before they were admitted to it. References to the sacraments and their significance are frequent in the writings of all the Fathers.\footnote{See contents and index references in “Ante-Nicene, Nicene, and Post-Nicene Fathers.”}

\randnotiz{Catechetical works in Middle Ages.} Catechetical works in the stricter sense we find first in the middle ages, but they are not yet called catechisms.

About the year 720 a monk of St. Gall in Switzerland, named Kero, is said to have written an explanation of the Apostles’ Creed and of the Lord’s Prayer in German. Achelis, “Practische Theologie,” page 102, disputes this and claims that Kero has been credited wrongly with what belongs to the following:

About 840 a monk in the cloister of Weissenburg, in Alsatia, named Ottfried, wrote a kind of Catechism explaining the mortal sins, the Apostles’ Creed, the Athanasian Creed, the Lord’s Prayer and the Gloria in Excelsis.\footnote{This book was published in Latin in 1713 by J. T. Eccard. It has always been highly esteemed.}

\quellseite{78}
In the eleventh century another monk, of St. Gall, named Notker, wrote a Catechism explaining the Apostles’ Creed, the Lord’s Prayer, the Song of Zacharias and the Magnificat.

A little later we find the first Catechism in the form of question and answer. It was written by Bruno, Bishop of Wurzburg, a man of royal blood and noble spirit. It explains in simple and clear questions and answers the Apostles’ Creed and the Lord’s Prayer.

In the fourteenth century we meet the great Jean Charlier Gerson, who has been called the evening star of the middle ages. Schutze, in his “Practical Catechetic,” calls Gerson the Spener of his time. He saw and mourned the failings and sins of the Church of his time. One of the finest intellects and probably the greatest scholar of his age, he yet delighted in teaching the young. He insisted that the main thing in all nurture is a living piety. He wanted a school in every parish, taught a school himself, and used the Socratic method in his instruction. He wrote “De Parvulis ad Christum Trahendis,” i. e., on Bringing Children to Christ, and “De Arte Audiendi Confessiones,” i. e., the Art of Hearing Confession. While he did not write a Catechism, he did insist on the instruction of the young when they came to confession. He also was a child of his time and opposed the giving of the Bible to the laity in the vernacular. \randnotiz{Gerson’s writings.}

\quellseite{79}
We come now to the pre-Reformation sects, the Waldenses, the Wiclifites, and the Hussites or Bohemian Brethren. \randnotiz{Waldensian Catechism.}


The Waldenses were evangelical Christians before the Reformation. They believed and taught the Formal Principle of the Reformation, viz., that all teaching must be in harmony with the Bible. They believed in the universal priesthood of believers. They were diligent and earnest readers and students of the Scriptures. They would receive no one into their communion who could not state and defend their teachings. Every household was a Christian school. They had a Catechism written in the form of questions and answers. It treated of the Commandments, Creed and Lord’s Prayer. A contemporary writes that it would have been hard to find a boy of twelve or fourteen who could not intelligently set forth and defend the evangelical teachings.

\randnotiz{Catechisms of Wiclif and Huss.} Wiclif, the great English character, aptly called the Morning Star of the Reformation, wrote a Catechism. Its contents were the Lord’s Prayer, Creed, Ten Commandments, thus reversing the Waldensian order. It has the strange title, “Pauper Rusticus.”


The great and good John Huss spent his last days before his martyrdom in writing a Catechism—contents, the Creed, the Decalog and the Lord’s Prayer. The Bohemian Brethren afterwards added questions
\quellseite{80}
and answers on Baptism and the Lord’s Supper, as also a Table of Duties made up of Bible verses.

The intrepid and illustrious Savanarola also wrote a “Catechism for Christian People,” containing expositions of the Commandments and of the Lord’s Prayer. \randnotiz{Of Savanarola}


When Luther was already upon the scene, but had not yet attracted wide attention, Nicolas Rusz wrote a Catechism with the obscure title, “A Threefold Cord” (Triplex Funiculus). On account of its evangelical contents the Romish authorities gathered and burned the copies of this book. \randnotiz{Of Rusz.}


From this brief survey we see that the Roman Catholic Catechisms of the Middle Ages had, as a rule, two divisions, viz., Lord’s Prayer and Creed, while the pre-Reformation sects’ Catechisms had three parts, viz., the Decalog, Creed and the Lord’s Prayer, though not all in the same order. The Bohemian Catechism as we saw was enlarged to five parts, with Table of Duties appended.

This brings us to the Age of the Reformation. \randnotiz{Early Reformation Catechisms.} We notice briefly the Catechisms that appeared before Luther’s.

In 1523 Urbanus Rhegius wrote an explanation of the Apostles’ Creed in twelve parts. The division was based on the ancient erroneous idea that each one of the twelve Apostles furnished a part of the Creed. A singular feature of this, the first evangelical
\quellseite{81}
catechetical book, is that the questions are asked by the pupil and answered by the teacher. A good hint for all Catechists.

Johannes Bugenhagen published a “Christenlehre” in 1524. In the same year Melanchthon wrote his “Handbüchlein: Wie man die Kindlein zu der Schrift und Lehre halten soll.” After an interesting introduction we find: 2, The Lord’s Prayer; 3, The Ave Maria; 4, The Sixty-sixth Psalm; 5, The Ten Commandments. This is followed by edifying expositions of a number of New Testament passages.

In 1528 Melanchthon followed this work with an expansion intended for preachers. This work he never finished. Another followed in 1549.

About 1524 Justus Jonas wrote: “Christlicher und Kurzer Unterricht von vergebung der Sünden und Seligkeit.” This is not, however, a Catechism, but a series of simple sermons for children.

With the exception of the first one mentioned, the above are all in the lecture form. None of them has the modern form of Catechisms.

\randnotiz{Brentz’s Catechism} Probably the most important work done before Luther was that by Johannes Brentz. He wrote “Fragestück des Christlichen Glaubens für die Jugend” in 1527. Nearly ten years later he thoroughly revised this book, using much of the material of Luther’s Catechisms. This Catechism is still held
\quellseite{82}
in high esteem and is widely used in Wurtemberg. In 1528 Andreas Althammer wrote “Katechismus, i. e., Unterricht Zum Christlichen Glauben.” About the same time Johannes Lachman wrote “Katechesis, i. e., Unterricht für Kinder: wie er in Heilbronn gelehrt und gehalten wird.” This was printed by Caspar Gratner, and is, therefore, often called the Caspar Gratner Catechism.

The last three were written in the form of question and answer. Other less important efforts were made by Johann Agricola, Conrad Sohm, Johannes Bader, Wolfgang Kopfel, Wolfgang Capito, Petrus Schultz, Christof Hegendorf, et al.\footnote{See Th. Harnack, “Der Kleine Catechisms Luther’s,” etc., 1856; G. Mönckeberg, “Die erste Ausgabe von Luther’s Kleinen Katechisms,” 1851. Von Zezschwitz Katechetik, pp. 224 ff.; G. Buchwald, “Die entstehung der Katechismen Luther’s und die Grundlage des Grossen Katechismus,” 1894. G. Kawerau, “Zwei Alsteste Katechismen der Lutherischen Reformation,” 1891, Achelis, Lehrbuch.}

\quellseite{83}
\chapter[Chapter VII. History of Luther’s Catechisms]{Chapter VII. The Catechist’s Text-Book. History of Luther’s Catechisms}

We come now to the text-book of the Lutheran Catechist. Luther has been called the Augustine of Protestantism. Along with his many varied and great gifts he had the donum docendi, i. e., the gift of teaching. In the midst of his herculean labors, he delighted in teaching the young. Perhaps no one before or after him had his qualifications for writing the Church’s Catechism. For simplicity, edification and devotion, for fullness and depth of contents, for correctness of arrangement, form and language, as well as for usableness in catechising, it has never been matched. \randnotiz{Matchlessness of Luther’s Catechism.} By its pure, powerful and spiritual portrayal of the essence of saving truth it became a mighty means for furthering the Reformation and has furnished the spiritual nerve and life-blood for the Mother Church of Protestantism from that day to this. Wherever the Church has learned, loved and lived its rich, yet simple lessons, she has prospered. Wherever she has neglected, vitiated or substituted its teachings, she has lost power and influence and life.

This booklet, so simple, so small, so beautiful, so
\quellseite{84}
full of heart and experience, cost Luther much arduous labor. As we read it, it all seems so natural that we might imagine that it flowed spontaneously out of the mind and heart of its great author. But not so. It has its history. It is something like a great sermon once preached by an earnest and able minister. He was asked how long it had taken to prepare that sermon. About fifteen years, was his answer. It was great, because it had in it the fruit of fifteen years of study, meditation and prayer.

Luther’s Catechism was a development, a growth, a fruit that had required years in the ripening. \randnotiz{Its cost to Luther.} Luther was a careful student of the Church, her history and her literature. He was at home in the Church Fathers. He believed in appropriating the good wherever he found it. He hesitated not to use Cyril of Jerusalem, or Cyprian, or Gregory of Nyssa, or Augustine, or Ottfried of Weissenburg, or the Bohemian Brethren, in his works. He knew the Catechisms of his friends and co-workers, enumerated above, and used them as he saw fit.

We note briefly the catechetical work that Luther did previously to writing his Catechisms.\footnote{“Practischen Theologie,” 1898, pp. 103–106. The essential parts of these pre-catechismal works of Luther have recently been made available for English readers in an inexpensive and indispensable book called “Luther’s Two Catechisms Explained by Himself,” by Dr. J. N. Lenker. The Luther Press, Minneapolis, 1908.}
\quellseite{85}
\randnotiz{Preceding catechetical work.} In 1517, during Lent, Luther preached on the Lord’s Prayer. He also lectured on the same subject to his students in the University of Wittenberg. He himself edited and published parts of this series in 1518, under the title “A German Explanation of the Lord’s Prayer for Simple Laymen, not for the learned.”

In the same year, 1518, he published, in Latin, an exposition of the Ten Commandments in popular sermons.

In 1520 he published his important preparatory work in German, entitled: “A Short Form of the Ten Commandments, The Creed and The Lord’s Prayer.”

This book is not a Catechism. It is rather an admonition to proper and intelligent confession for those who purpose to go to the Lord’s Supper. In 1522 Luther revised and enlarged this work and added the Ave Mariam and Meditations on several Psalms.

In his preface to the first edition of this book Luther says: “It has not come to pass without the special ordering of God that the common Christian who cannot read and understand the Scriptures can learn and know the Ten Commandments, the Creed and the Lord’s Prayer, in which three parts everything is contained essentially and fully which it is necessary for a Christian to know.”

\quellseite{86}
\randnotiz{Luther’s booklet for the laity and children.} Perhaps nearest to a Catechism is the little book of 1525, called “A Booklet for the Laity and Children.” This contains the A, B, C, which distinguishes between vowels and consonants, with passages from the Sermon on the Mount, the Lord’s Prayer, the Institution of Baptism, and the Lord’s Supper. While we have here, therefore, a book of elementary and religious instruction for the young, it is by no means a Catechism. After writing it Luther wrote to his friend, Nikolaus Hansmann: “The Catechism I postpone, as I wish to do this work at one stretch.”

In 1526 the first edition of his “German Mass and Order for Divine Service” appeared. In the Introduction to this he also writes: “In the German Church Service a simple Catechism is sorely needed.” Luther felt that the people cannot be expected to worship correctly, or to get the full benefit of the sermon unless they are informed on the fundamental teachings of the divine Word. He felt that the sermon needs the Catechism; the preacher, the Catechist. Neither function is complete without the other. They complement each other.

The immediate occasion for the writing of the Catechisms was the great need of the Church, as it was brought so vividly before him in the visitation of the Churches in 1529. This need is portrayed in an impressive manner and in realistic colors in his \randnotiz{Need, the occasion for the Catechism.}
\quellseite{87}
report of that visitation. As the essential part of that report is contained in his Introduction to his Catechism, and as that Introduction is found in most good editions of the Small Catechism and in The Church Book we will not quote it here. It goes without saying that every teacher of Luther’s Catechism will make himself familiar with that ringing Introduction, whose admonitions are sorely needed even to-day.

This brings us to the story of the writing of the two Catechisms.

\randnotiz{Which first, Large or Small Catechism?} Much controversy has taken place as to which one of the two, the Large or the Small Catechism, was written first. The commonly accepted view has been that Luther intended to write but one, and started with that intention. But that as he wrote, the work grew in his hands until it became too bulky for family or school use, and that for this reason, after he had finished the large, he started anew and wrote the small one.

Through the investigations of historical criticism this view will have to be given up. The whole question is discussed in a thorough and scholarly manner by Achelis in his “Practische Theologie.”\footnote{See pages 31, 32.} Achelis informs us that G. Buchwald had made several of those surprising documentary discoveries that have played so large a part in recent historical
\quellseite{88}
and theological discussions.\footnote{See “Die Entstehung Der Katechismen Luthers,” p. xi. ff.} \randnotiz{Buchwald’s discoveries.} This Buchwald found in the Law-School Library in Zwickau, a collection of \emph{Lutherana Inedita} as well as a number of heretofore unknown letters of a certain Röre. In 1893 he also found in the library of the University of Jena thirty-three volumes of Luther’s sermons. From these important finds Buchwald concludes that Luther began his work with the purpose to furnish the official Church visitors and both the preachers and sextons (Küstern), of the parishes visited with a handbook for carrying out the duties laid upon them. The visiting preachers were to order the parish pastors to preach regularly on the five parts that Luther had selected and that were to become the divisions of his Catechisms. The sextons were to instruct or catechise the youth of the parish in these five parts. But there was no suitable book available to assist the inexperienced to carry out these duties.\footnote{G. Kawerau, “Zwei Alteste Katechismen,” p. 3 ff.} Luther also intended to furnish the house-fathers with a guide for the instruction of their families. But while Luther was thus working out a handbook for pastors and sextons, by working over the sermons that he had preached on the five parts, he soon found that this would not be a suitable book for the plain heads of families and their children. He therefore stopped in the midst of the
\quellseite{89}
work that grew into the Large Catechism and wrote his explanation of the Commandments, the Creed and the Lord’s Prayer, as we now have them in the Small Catechism.

These three parts he had printed on large wall tablets and sent out January 20th, 1529. On March 16th he sent out another set of tablets, containing the explanation of Baptism and of the Lord’s Supper. It appears, therefore, that while Luther started to write the Large Catechism first he broke off this writing and wrote the Small Catechism before he took up a second time and finished writing the large one. The Small Catechism was, however, first published in the form of wall tablets in separate parts. This was done before the writing of the Large Catechism was finished. In the form of tablets, therefore, the Small Catechism has priority in time. \randnotiz{Priority of Small Catechism.} In book form, however, it did not appear until the large one was on the market. It was for the book form of the Small Catechism that Luther wrote his Preface. It therefore need not embarrass us that this Preface refers to the Large Catechism, which had now been finished and published. The first edition, in book form, with Preface, appeared June 13th, 1529, and bears the title: “Enchiridion: The Small Catechism, for the Common Pastors and Preachers, Martin Luther, 1529.”


After Luther had written his little book he loved
\quellseite{90}
it. He kept on revising, improving and enlarging it. The first Wittenberg edition of the book contained,\randnotiz{Additions to first edition.} in addition to the five parts of the tablets, Luther’s Preface, Morning and Evening Prayers, A Table of Duties, and a Marriage Formula.\footnote{This first book-edition appeared May 16th, 1529.} To the third edition,\footnote{This edition was illustrated with twenty crude Biblical pictures.} of June 13th, he added his \emph{Taufbüchlein}, Directions for Confession and the Litany. In the edition of 1531 he left out the direction of Confession and put in its place the three questions and answers on Confession, and himself placed them between the fourth and the fifth parts. In this edition he also left out the Litany and added the Introduction to the Lord’s Prayer. The form of this Introduction was afterwards modified and received its present wording in 1534. In this year he also made some changes in the proof-passages of the Table of Duties.


The part on “The Power of the Keys,” which some Catechisms call the sixth part, was never written by Luther and was not added to the Catechism until after Luther’s death.\footnote{The essence of this part was taken from the Sermons for Children, which Brenz prepared for Brandenburg-Nürnberg K. O., in 1533.} \randnotiz{“The Power of the Keys.”}


Likewise the “questions for those who intend to go to the Lord’s Supper,” found in some Catechisms, though based on a Maunday Thursday Sermon
\quellseite{91}
of Luther’s, were prepared by Johann Lang, in Erfurt, and added to the Catechism after Luther’s death.\footnote{It is significant that Luther always speaks of the five parts of his Catechism, and never, as far as we know, of six.}

\randnotiz{Sources} It remains for us to notice briefly the sources of Luther’s Small Catechism. We have seen that Luther was Catholic in his knowledge and appreciation of truth. He prized it wherever he found it, and used it wherever it suited him. He hesitated not to borrow, but he did not borrow all that some writers claimed he did.

It has been claimed by eminent writers\footnote{Dr. C. P. Krauth, “Conservative Reformation.”} that he borrowed from Kero, of St. Gall, in his explanation of the Apostles’ Creed. Achelis, as we saw above, denies this and claims that Ottfried’s Weissenburg Catechism contains the passages said to be borrowed and that Kero never wrote them. That if Luther borrowed, he borrowed from Ottfried, of which, however, there is no proof. That there is a striking similarity between Luther and Ottfried and Cyprian on the Lord’s Prayer is shown by Achelis in parallel columns.\footnote{“Praktische Theologie,” p. 113.} He explains, however, that both Luther and Ottfried probably borrowed from Cyprian, and that Luther never saw Ottfried’s book.

Be this as it may, we do know that Luther
\quellseite{92}
borrowed from Cyril and from Cyprian and from other Church Fathers.\footnote{See Schütze, “Praktische Katechetik,” pp. 52, 53, 60, 61.}

Luther also had acquaintance with the Catechisms of the Waldenses and the Bohemian Brethren.\footnote{Dickhoff, “Die Waldenser im Mittelalter”; Herzog, “Die Romanischen Waldenser”; Zezschwitz, “Die Katechismen der Waldenser und Böhmischen Brüder.”}

The Catechism of the Bohemian Brethren was sent to Luther by Bishop Lukas. It was in Latin, and the bishop begged Luther to translate it into German and publish it for Germany. Luther examined it carefully, found much to commend in it, but could not agree with its teaching on the Lord’s Supper. He therefore declined to translate it.

Naturally Luther also knew the Catechisms written before his. He thought highly of the two Catechisms of Brentz. When Opsopius translated Luther’s Large Catechism into Latin, he appended the Brentz Catechisms. Palmer says: “The story of the origin and development of Luther’s Catechism is an important chapter in Church history.”

It certainly was no private writing for private use. Luther intended it for the establishment and security of the Church. \randnotiz{Luther’s intention.} It was to be the instruction book for the Church. He says of it: “This teaching and faith I cannot arrange and set forth better than in the order and form in which it has been done
\quellseite{93}
from the beginning of Christianity, is still done and so remains to-day; that is in the three parts, the Ten Commandments, the Creed and the Lord’s Prayer. When these three parts have been appropriated, one must also be able to answer as to the Sacraments which Christ Himself ordained, \emph{i.\,e.}, Baptism and the Sacrament of the Sacred Body and Blood, as is related by Matthew and Mark at the end of their Gospels. Thus we have for all time, in five parts, the whole Christian doctrine which shall be taught everywhere.” From the story of the composition of the Lutheran Catechist’s textbook we see that its author was specially gifted and fitted for his work, and that in this little book he gave the Church his very best and most permanent gift. Of this more in the next chapter.

\quellseite{94}
\chapter[Chapter VIII. Characteristics and Value]{Chapter VIII. Characteristics and Value of the Catechist’s Text-Book}

\randnotiz{Luther’s estimate.} Luther himself prized his Catechism highly. He says: “This Catechism is truly the Bible of the laity, wherein is contained the entire doctrine necessary to be known by every Christian for salvation. Here we have:

“First, the Ten Commandments of God, the doctrine of doctrines by which the will of God is known, what God would have us do, and what is wanting in us.

“Secondly, the Apostles’ Creed, the history of histories, or the highest history, wherein are delivered to us the wonderful works of God from the beginning, how we and all creatures are created by God, how all are redeemed by the Son of God, how we are also received and sanctified by the Holy Ghost and collected together to be a people of God, and have the remission of sins and everlasting salvation.

“Thirdly, the Lord’s Prayer, the prayer of prayers, the highest prayer which the highest Master taught, wherein are included all temporal and spiritual blessings and the strongest comforts in all temptations and troubles and in the hour of death.
\quellseite{95}
“Fourthly, the blessed Sacraments, the ceremonies of ceremonies, which God Himself has instituted and ordained, and therein assured us of His grace.”

Matthesius, a contemporary of Luther, writes:

“If Dr. Luther in his whole life had done nothing good except to have brought his two Catechisms into the homes, the schools and the churches, the world could never thank him enough.”

\randnotiz{Estimate of others.} Justus Jonas says:

“The Catechism is a little book. It can be bought for sixpence, but six thousand worlds could not pay for it. If the Lutheran Church had brought no other benefit into the world than that she made this Catechism known to the people, she would have done more to build the Church than all the higher schools of learning that have been on the earth.”

The celebrated theologian, Dr. Lyserius, writes:

“It can be said with truth that in this little book there is such a depth and overflowing of divine teaching that if all evangelical preachers would preach during their whole life this hidden wisdom of God to the common people, as it is here given in few and simple words, they would not exhaust its depths. Therefore let no one think himself too high to exercise himself in the Catechism.”

John Arndt, in a sermon on the Catechism, says:

“The Catechism is a brief instruction in the Christian religion, and includes in itself the doctrine of
\quellseite{96}
the law of God, Christian faith, the Lord’s Prayer, the institution of Holy Baptism and of the Lord’s Supper, which five parts are an epitome and kernel of the entire Holy Scriptures, for which reason it is called a ‘little Bible.’”

Dr. Seiss, in his “Ecclesia Lutherana,” says: “It is the completest summary of the contents of the Bible ever given in the same number of words. It gave to the reviving Church a text-book for the presentation of the truth as it is in Jesus to the school, class-room and Church.”

Dr. Krauth says: “The Catechism is a thread through the labyrinth of divine wonders. Persons often get confused, but if they will hold on to this Catechism it will lead them through without being lost. It is often called the ‘little Bible’ and the ‘Bible of the laity,’ because it presents the plain and simple doctrines of the Holy Book in its own words. Pearls strung are easily carried, unstrung they are easily lost. The Catechism is a string of Bible pearls. The order of arrangement is the Historical, the Law, Faith, Prayer, Sacrament of Baptism and all crowned with the Lord’s Supper—just as God worked them out and fixed them in history.”

Thus we might go on quoting page after page of words of admiration and praise from the greatest minds in our own and in other churches. And we are not prepared to charge these men with extravagance
\quellseite{97}
in their utterances. Indeed, the more we examine and study the contents of this little book, the more we are convinced that it is unique and admirable in matter and in plan. And yet it has been criticised both as to its matter and its plan.

\randnotiz{Published in many languages.} Before we notice some of these criticisms, it may be well to be reminded that by the end of the sixteenth century it had been published in Latin, Greek, Hebrew, Spanish, Dutch, English, Danish, Norwegian, Swedish, Icelandic, Finnish, Lithuanian, Esthonian, Polish, Bohemian, Croatian, together with numerous dialects.

We note now a few of the principal criticisms of Luther’s Small Catechism. It has been criticised both as to its contents and as to its arrangement.

\randnotiz{Dinter’s criticism.} Perhaps the most important critic is Dinter, of whom Schütze says:\footnote{“Praktische Katechetik,” p. 69.} “He is, without doubt, the greatest master of catechetical art.” In his book, “Die vorzüglichsten Regeln der Katechetik,” Dinter says: “With the deepest reverence for Luther’s merit, it would yet be hard to deny that his Catechism is marred by misproportion, incompleteness, lack of order and lack of clearness.” He claims in general that Luther’s Catechism is defective in that it fails to give the doctrine and proof of the being or existence and nature of God, of His perfections and attributes, of the divinity of Christ and of the
\quellseite{98}
\randnotiz{Scriptures.} Scriptures. It fails also in ethics, in that it does not set forth and enforce the duty of caring for one’s own welfare nor the duty of caring for the souls of our neighbors. This is the essence of Dinter’s criticism.

It is not needful that we take much time in answering this criticism. Suffice it to say:

First. Our book is a Catechism and not a treatise on systematic theology. It is an instruction book for children and young people, the baptized youth from Christian families.

Second. These young disciples are neither heathen nor atheist. Every preacher is sent forth not to prove, but to preach the Gospel. Even more so, \randnotiz{Answer to Dinter.} the Catechist is not to prove, but to teach the elements of the Christian faith. It is enough that the Commandments begin with “I am the Lord thy God,” and the Creed with “I believe in God,” and the Lord’s Prayer with “Our Father who art in heaven,” and that the Sacraments exist because of the Lord’s words of institution. Here is ample room and encouragement for the wise Catechist to enlarge and deduce and exemplify with facts and fortify the existing belief, and show the attributes from the words and works of God.

Well does Steinmeier say: “A system of theology at the threshold of the Catechism would be at the wrong place. It would anticipate what should follow
\quellseite{99}
and develop as a matter of course. It would be a pedagogical error and would be violence to the psychology of the child.”

The ethical objection of Dinter seems even more groundless. What is the whole Catechism but an admonition and explanation as to the care of one’s own soul? And what is the second table of the Law, with Luther’s masterful explanations, but an earnest admonition to have a constant care for the body, soul, reputation and estate of our neighbor? Surely the Catechist who knows how to bring out and apply what is in Luther’s Catechism, who knows how to make his own teaching clear and drive it home, will find no important fact, doctrine or duty missing.

\randnotiz{Kraushold’s criticism.} A criticism of a different kind in Kraushold’s “Katechetik,” page 29, where he affirms that the teaching of Luther’s Catechism in the Reformation era led to an over-emphasizing of the true doctrine, while that which ought to be the real aim of all instruction, the inner apprehension of the truth and the development of the subjective life of the youth, was crowded into the background.

That such has often been the case, and that such superficial lifeless and loveless teaching is the bane of much of the catechising of to-day is only too true. Of this more hereafter. But wherever it is and has been the case, the fault is certainly not in
\quellseite{100}
the text-book. We have often said that Luther’s Small Catechism is pre-eminently a book of personal and experimental religion. It is surprisingly sad that this is so little recognized and emphasized. \randnotiz{Answer to Kraushold.} Take the explanations of the Commandments. “What is meant by this Commandment?” The answer is not an objective, theoretical setting forth of the deep and divine import, but, in every Commandment “We should fear, love and trust.” Fear and love and trust are experiences. Each and all are personal, vital. Look for a moment at the Creed. Note the “I” and the “me” and the “my” as they form the basis of all three explanations. Was there ever a more beautiful, simple, reverend, devout and tender expression of heartfelt experience, apprehension and appreciation than is expressed in the explanation of the second article? Or of the third article? Truly here is the expression and confession of a trusting, clinging, adoring, living and loving faith that pulses warm with the life-blood of a healthy inner spirituality. Truly here we have a Catechism that can be prayed from beginning to end.

Even more outspoken have been the different views of the friends of the Catechism as to the wisdom of the arrangement, and as to Luther’s purpose in this arrangement.

Did Luther intend that the first part, \emph{i.\,e.}, the
\quellseite{101}
\randnotiz{Questions of friends.} Law, should be used pedagogically, to bring about a knowledge and sense of sin, guilt and helplessness, that would drive the pupil to seek release and refuge in Christ? Or did he intend that, for the baptized children of the Church, the didactic, or the third use of the Law, should be inculcated?

Did Luther intend that his Catechism should form a doctrinal system or plan of salvation, or did he intend that each part should be taught and impressed in all its bearings, without regard to its relation to the other parts? These are the questions on which Catechists have divided.

\randnotiz{Argument for pedagogical use of Law.} We consider, first, the question as to the use of the Law. Zezschwitz is the most prominent and the ablest of those who believe that the Law in the Catechism has the same office and purpose that it has in the general order of salvation, \emph{viz.}, to awaken a sense of guilt and fear, and so to lead to Christ. He bases this view, as do others also, on an extract from Luther himself, where he says:\footnote{“Kurze Form der Zehn Gebote, des Glaubens und des Vater Unsers,” 1520.} “Just as a sick one needs to know, first of all, what his sickness is, and after that to know what the remedy is that will help him, and thirdly, to earnestly desire that remedy and to ask for it and have it brought to him.”

But Zezschwitz forgets that Luther wrote this in
\quellseite{102}
1520; that he was not writing a Catechism, but a brief instruction for confession; that this was not for Christian children, but for adults who do not know the evangelical teaching concerning confession and absolution. For such, certainly, the law is to awaken a sense of guilt and fear, and so a readiness to fly to the Saviour as He is set forth in the Creed, and to cry to God the Father, for the sake of Christ the Son, as he learns in the Lord’s Prayer.

But did Luther mean the same thing by placing the Commandments before the Creed and by his explanation of them in the Catechism? Achelis denies this most positively and quotes Steinmeier on the Decalog, 1873, Leopold Schultze,\footnote{“Bausteine zur Erklärung von Luther’s Katechismus,” 1891.} J. Gottschick,\footnote{“Luther als Katechet,” 1883.} K. Knoke,\footnote{Luther’s “Kleiner Katechismus,” 1886.} Rhoden,\footnote{“Zur Katechismusfrage,” 1890.} Kolde,\footnote{“Der Kleine Katechismus,” U.\,S.\,W., 1895.} B. Dörries,\footnote{“Der Glaube,” 1895.} Th. Hardeland, \emph{et al.}\footnote{“Der Kleine Katechismus nach Luther’s Schriften,” 1889.} To these we may add the American Professor M. Reu, of the Theological Seminary of the German Iowa Synod at Dubuque, who has written what is probably the most exhaustive and scholarly work on the history of Lutheran catechisation ever before attempted.\footnote{For his view on the question before us, see his Article in \emph{Lutheran Church Review}, July, 1905, also his introduction to his Luther’s “Small Catechism with Explanations,” English, Wartburg Pub. House, Chicago.}

\quellseite{103}
\randnotiz{Argument for didactic use of Law.} All these authorities differ from Zezschwitz and maintain that Luther intended to use the Law for the baptized children of the Church, mainly didactically, \emph{i.\,e.}, to show them what their dear Father desires of His loving children, without overlooking the pedagogic use which shows the remaining sin, and should keep awake that daily repentance which all the children of God constantly need. With this view, after a careful study of both sides of the question, we fully agree.

Turning now to the other question, as to the arrangement and order of the five parts, we find that it is closely related to the former one and is already partly answered.

\randnotiz{Argument for Catechism as a doctrinal system.} Zezschwitz maintains, with many others, that there is a vital connection and relation between the parts, and that no single part can be rightly understood where its relation to the other parts is overlooked. These writers claim that the order is historical as well as dogmatical. That just as in the history of God’s revelation, we have first Moses, then Christ, then the Holy Spirit, then the ability to pray, and then the gift of the Sacraments; so we have it in the Christian life, and so we have it in the Catechism. Many fine and fanciful theories and speculations have been written along this line.

\quellseite{104}
Achelis\footnote{“Praktische Theologie,” p. 149.} grows indignant and sarcastic in combating this view. He claims that it is “religious Darwinism.” That it was Darwin’s great discovery that the stages of development in the evolution of each genus are repeated in the history of each individual of the species, and that the epochs in the evolution of the race are repeated in the life of each individual. To claim, therefore, that the history of God’s people, who had first law, then the Spirit, is repeated in each Christian life, is to transfer the fancies of natural evolution into the sphere of Christian life and Christian teaching.

This is at least interesting, if nothing more. It does not answer our question.

\randnotiz{Answer.} Achelis and all the eminent Catechists quoted against Zezschwitz on the former question, believe that Luther did not intend that his Catechism should be an organic system of doctrine. They believe that it was Luther’s idea that each part should be taught by itself, without regard to a supposed relation to the other parts. This is brought out clearly and forcibly in the Church Review article of Professor Reu, referred to above.

All these writers believe that each one of the five parts contains the essence of Christian Truth: That if the Catechist is master of his art, he will know how to draw it all out of the part and Luther’s explanation
\quellseite{105}
of it. That Professor Reu knows how to do this he shows in his own explanation of the Catechism.

Now our own conviction in brief is this:

\randnotiz{Author’s conclusions.} Too much has doubtless been made of the supposed system. Luther himself certainly did not lay stress on it. In his Introduction to the Catechism he complains three times over of the people that they did not know “either the Lord’s Prayer, or the Creed, or the Commandments.” This is certainly not the order of parts in the Catechism. In his warning to the Christians in Frankfurt against the Zwinglian teaching he admonishes the pastors to examine those who come to confession, “whether they know the Lord’s Prayer, the Creed and the Commandments.” This was written in 1533, and is not the order of the Catechism. He tells these pastors that in such an examination they will “see that the Creed, the Lord’s Prayer and the Commandments comprehend a short form of teaching for the youth and simple people, and have been called Catechism from the beginning.” We might give other quotations, but these suffice to show that Luther laid no emphasis on the order, and therefore did not profess to build a system.

And yet, while convinced that there is no dogmatic system, and that the Catechist should not worry himself in building labored and fanciful
\quellseite{106}
\randnotiz{Order natural and beautiful.} bridges between the parts, we believe that there is an instructive naturalness and beauty in the order of the parts.

We can, therefore, still conclude our characterization and appreciation of Luther’s Small Catechism with the words that we wrote nearly a quarter of a century ago:

“I come into this world ignorant, yet full of presentiments and questions. I learn my first vague lessons about myself and God. I naturally ask: For what purpose has God put me here? What does He want me to do? The Catechism answers: To do His will, to keep His commandments. Here they are, and this is what they mean. I study them, and the more I study the more am I convinced that I never did and never can keep this law perfectly.

“I ask again, What shall I do? My Catechism tells me that I must have faith. I must believe. But what shall I believe? Answer, this summary of truth called the Apostles’ Creed. It tells me of my Creator, His work and providence, and His gift of a Redeemer. It tells me of that Redeemer and His redemption, of the gift of the Spirit and His application of redemption. It not only tells me what to believe, but in the very telling it offers me strength to believe.

“But I am still weak and more or less perplexed. Whither shall I go for strength and grace? My
\quellseite{107}
Catechism gives answer, Go to the great Triune God. Ask Him in prayer. Here is a model. It will teach you how to pray. I learn what it is to pray. But again I ask, How do I know that God will hear my prayer? Is He interested in me personally? Has He any other means besides His written Word to assure me of His love and to give me, in answer to my prayer, strength to believe Him and love Him? My Catechism points me to my baptism. It teaches me what it means, and how that in it I have God’s own pledge that He is my Father and that I am His child. Here, then, is a fountain to which I can return again and again when weak and perplexed.

“Further, my Catechism teaches me concerning my Saviour’s last legacy of love before His death for me, His Holy Supper. In it He holds out to me and gives to me, personally and individually, Himself and all His heavenly grace.

“Thus does this little Catechism meet me in my perplexity, take me by the hand and lead me through the labyrinth of the wonders of grace. Thus does it tell me what I am, what I need and where and how to get what I need. It takes me to the wells of salvation. It draws from them living water. It holds it to my thirsty lips. It gathers the precious manna of the word and feeds me when I am faint and weary.”\footnote{“Way of Salvation in the Lutheran Church,” pp. 79, 80.}
\quellseite{108}
\chapter[Chapter IX. The Question and Answer Form]{Chapter IX. Characteristics and Value (Continued). The Question and Answer Form}

There have been objections against Luther’s question and answer method.

It is objected that it is unhistorical. The early Christian Catechists, it is claimed, did not teach their pupils by question and answer. \randnotiz{Objections to question and answer method.} They employed the usual method of setting forth truth. This method was continuous discourse. They put their instruction in the form of theses, propositions, positive statements. These were then elaborated, illustrated, discussed, applied, and, if need be, defended. The form was thetical and not interrogative, acroamatic and not erotematic.

As to this historical claim, we freely admit that there is truth in it. But it is not the whole truth. \randnotiz{Antiquity of questioning method.} The questioning method has nearly always been used interchangeably with the thetical. Some of the greatest teachers of the ages made much use of the question method. Pythagoras, who lived five hundred years before Christ, and was the greatest teacher of his day, was largely erotematic in his method. The greater Socrates, one hundred years later, became the master and the model for all ages

\quellseite{109}
in this method. The greatest of all teachers, Jesus Christ, used this method again and again. So did the Apostles, especially the greatest of them, the inspired dialectician, Paul. True, Cyril of Jerusalem gave his instruction in the form of popular sermons. Thirty-three of these catechetical discourses are extant. They deal with expositions and applications of lessons from the Gospels, the Apostles’ Creed, the so-called mysteries, viz., Baptism, Confirmation and the Lord’s Supper.

The great catechetical discourses of Gregory of Nyssa are similar in form.

Augustin wrote a work for Catechists, “De Catechizandis Rudibus.” In this work he shows the use of proper questions and urges their frequent use.

About the only form of catechising that was left in the middle ages was the examining of the communicants in the private confessional. Here the priest was expected to examine and instruct both as to the understanding and the life of the confessing one. This was always done in the erotematic form.

\randnotiz{Difference in pupils.} Coming back again to the great Catechists of the early Church, we must always bear in mind that their pupils were not children, as ours are, but mature men and women, seekers after truth from among the Jews and the Gentiles, often the keenest intellects and best informed of their day. With such a class the instruction would of necessity take

paragraphannotation?
\quellseite{110}
a more formal character, be continuous, expository, apologetic and argumentative. Here the eloquence and the persuasiveness of the orator and the advocate would naturally find place. And yet, even here, there were frequent changes to colloquy, to question and answer and free discussion.

Another objection to Luther’s question and answer method is that it prevents a clear and connected setting forth of the truth to be taught; that the necessary breaking up of the matter into fragments destroys its effectiveness and robs it of its power.

To this we answer that it all depends on the nature of the questions and answers. \randnotiz{Effectiveness of Luther’s questions and answers.} The objection certainly does not hold against Luther. His questions are either a part of the continued elucidation of the truth, or they emphasize particular points of the truth; or again, they prepare the mind for a clear conception of the point to be brought out next. They are calculated, as all good questions should be, to awaken interest, to arouse attention, to prepare the mind for the grasping and holding of the truth contained in the answer and to make it as impressive as possible. What Catechist or catechumen has ever felt that the body of truth has lost force or meaning by Luther’s questions and answers? The objection is imaginary. If it holds good against any explanation of Luther’s Catechism,
\quellseite{111}
then the questions and answers of that explanation are faulty.

We are by no means convinced, therefore, by the objections of the able but rationalistic Dinter against Luther’s method, nor by the recent arguments of the Rev. Dr. Joseph Stump, in the Introduction to his most excellent, full and valuable explanation of Luther’s Small Catechism. Neither does the thetical explanation itself convince us. \randnotiz{Thetical method not satisfactory.} The author himself seems to feel the need of questions and introduces a series of his own after the explanation of each paragraph or section.

Although Professor Reu is more consistent and does not introduce questions into his thetical explanation of Luther, our objection holds against his unique, suggestive and highly esteemed work also. For children’s classes we still believe in questions and answers. Did we have all ideal children, then we might try more ideal methods. But since we must still deal with the real class, and since we always have the tryingly dull and lazy, the vexingly trifling and mischievous, the harassingly heedless and hostile, we still need and are glad we can avail ourselves of the tried, tested and ever helpful method of question and answer in catechising.

But this does not mean that we would discourage the more lengthy and unbroken explanation, defence
\quellseite{112}
or application. No true Catechist can do without these. They will come in again and again. The questions and answers will lead up to and demand them. \randnotiz{Both methods used by a good Catechist.} The right catechetical instinct will direct as to their place and length. It will prevent lecturing and preaching to the class. But it will give rise to earnest, effective and never-to-be-forgotten explanations and admonitions in every lesson. A teacher who confines himself to the questions and answers of the book is unfit to be called a Catechist. There is something seriously wrong in his spirit or his ability. We believe most heartily in a combination of the acroamatic and the erotematic methods.

The good Catechist will never confine himself to the questions of the book, however good they may be. He will always interject questions of his own. \randnotiz{Free questions.} To these the pupil must answer in his own words. In this manner he will test the comprehension of the pupil. He will know whether the pupil sees the bearing of the matter on his own heart and life. The pupil needs to know two things on each distinctive point brought before him. He needs to know first, what does it mean in itself? and second, what does it mean to me? All this can be brought out only by direct, free questions. The teacher’s questions ought to arouse interest and inquiry in the pupil. He will use his best endeavors
\quellseite{113}
to call out questions in the pupil. He is doing his best catechising when the pupil asks more questions than the teacher, provided always that the questions are along the line of the subject in hand. A good catechetical hour is one in which there is an opening of heart to heart, free inquiry, free discussion, asking and receiving light, help, hope, encouragement, determination.

\randnotiz{Art of right questioning.} The art of right questioning is no small part of the good Catechist’s equipment. He must learn how to frame his questions. Like Luther’s, they must be made up of simple words. They must be short. They must form a simple sentence. Complex or compound or mixed sentences of any kind are inadmissible in catechisation. The questions must be apposite, i. e., they must fit the answer sought. They dare not be ambiguous; must be capable of only one meaning and of one answer. A good catechetical question should be so framed that the pupil can find the answer in what he has already learned, or in his own experience, or in his own understanding. Sometimes it is allowable to ask a question to which it is impossible for the pupil to find an answer. The purpose, then, is to awaken a desire for an answer and to prepare the pupil for the reception of the teacher’s answer. It goes without saying that the teacher must really answer the question, so that the child can grasp it and rest in it.
\quellseite{114}
All empty, useless or would-be smart questions are to be scrupulously avoided by the teacher and repressed in the pupil.

The so-called yes and no questions are to be avoided. One may be permitted only when it is needed as a stepping stone to the real question. They are calculated to encourage laziness and guessing. And yet the otherwise good old Catechist, Rambach, uses them and advises them.

\randnotiz{Questions with part of answer.} Questions which contain a part of the answer are also to be used sparingly. They may be allowed to encourage a timid pupil, but not for a lazy pupil. We believe that Luther’s questions come up to all these requisites of good questions. In proportion as an explanation of Luther’s Catechism meets these requisites in that proportion it is a good explanation. And in proportion as a Catechist has the ability to frame his free questions after these requirements, in that proportion is he, on this point, a good Catechist. It is understood, as a matter of course, that the answers are to be as good as the questions. To this whole subject we shall return later.

\quellseite{115}
\chapter[Chapter X. The Story of Its Use]{Chapter X. The Catechist’s Text-Book. The Story of Its Use}

We have seen that Luther’s Small Catechism is most admirably adapted to its purpose. We have noted a few of the many tributes that have been given to it. We might have filled pages with these tributes. We have noted some of the criticisms and objections that have been made. We have seen that they are really groundless. The superior value of our text-book is still intact.

But after all, the proof of the pudding is in the eating. Luther’s Catechism is nearly four hundred years old. It has had ample time to be tested. \randnotiz{The test of use.} Has it stood the test of use?

It leaped into prominence as soon as it was published. \randnotiz{Rapid and widespread adoption.} Its circulation reached over one hundred thousand in the generation to which it was first given. In that century it was translated and circulated in Latin, Greek, Hebrew, Arabic, Syriac, Spanish, French, Dutch, Danish, Swedish, Norwegian, Icelandic, Finnish, Lithuanian, Esthonian, Bohemian, Polish and Croatian. In the Lutheran schools of Germany and other lands it became a
\quellseite{116}
regular text-book, and soon crowded out all others.\footnote{Book of Concord, Vol. II., p. 50.} As Dr. Jacobs says: “It has made its way into almost every living language of the world. It bears its testimony to Christ in every part of the earth, from Greenland and Iceland to India and Southern Africa, in the church, the school, the family, equally at home in the palaces of royalty and in the huts of poverty.”\footnote{Jacobs’ “Lutheran Movement in England,” p. 104 ff.}

Eight years after the publication in Germany a group of devout, earnest and scholarly Englishmen were gathered together in Archbishop Cranmer’s house to draw up a book that was to set forth in simple form what the people were to believe. The result of several months of conference was the publication of a book called “The Bishop’s Book,” or “Institution of a Christian Man.” This most interesting and devout little book, by which the English Reformation reached the highest point attained in the reign of Henry VIII., was largely a paraphrase of parts of Luther’s Catechism.\footnote{Ibid.} In 1548 Cranmer published his “Catechismus.” This was based on sermons on Luther’s Catechism.\footnote{Ibid., pp. 314–322.} \randnotiz{English Protestantism’s indebtedness.} Every Lutheran ought to read the most interesting story of what English Protestantism owes to Luther. And no Lutheran Catechist can afford to be ignorant
\quellseite{117}
of the influence of Luther’s Catechism in England. Get and read Dr. Jacobs’ “Lutheran Movement in England.”

\randnotiz{Catechism in America.} Luther’s Small Catechism was the first book ever translated into the language of the American Indians,\footnote{Read the fascinating story in Dr. Jacobs’ “History of the Evangelical Lutheran Church in the United States,” p. 82 ff.} by John Campanius, a Swedish Lutheran, in 1646. The first edition in America was published by Count Zinzendorf in 1744. Five years later an edition approved by Lutheran pastors appeared from the press of Benjamin Franklin. In the same year the first English translation in America appeared. Since then hundreds of editions have appeared in America, in nearly all languages spoken in our land.\footnote{See two interesting articles by Dr. B. M. Schmucker in Lutheran Church Review, April and July, 1886.}

To return to Germany and other European countries. Space forbids our going into a detailed account of the numerous explanations and expansions of Luther’s five parts that were worked out, circulated and used. \randnotiz{First hundred years in Germany.} For a hundred years after its first appearance, instruction consisted mainly in having the children memorize Luther’s Catechism. Further explanations were left to sermons on the Catechism, which grew into a distinct and large species of sermonic literature. They were generally modeled after Luther’s Large Catechism, and
\quellseite{118}
the “Nuremberg Kinderpredigten” of 1533. The children were often questioned on these sermons. In this way a mere mechanical memorizing was guarded against as though it were in itself sufficient. Earnest endeavors were made to have the catechumens understand what they had memorized.

\randnotiz{Spener’s cautions and influence.} To this end Pietism was especially zealous. Spener, in the Introduction to his “Erklärung der Christlichen Lehre nach der ordnung des Kleinen Katechismus, D. Martin Luther’s,” says: “As to the use I would have made of this work, my opinion is not that the main thing is the memorizing of the words of the answers. Against a mere mechanical memorizing I would rather utter a warning than say aught in its commendation. I am much more concerned that by reading, examining and discussing of the questions the meaning shall be made so clear and impressive, that the catechumens and all students of this booklet shall be able to reproduce the meaning of the matter in question out of his own understanding and in his own words. . . . And further, the learners should not only know what they ought to believe, and understand the same, but should become accustomed to ask what principle or rule of life, what help and encouragement to personal godliness is contained in each particular truth.”

This Catechism of Spener’s, from whose Introduction
\quellseite{119}
this quotation is given, is one of the most useful we know for the Catechist. Every Catechist who will earnestly and regularly study Spener’s Explanation before he goes to his class will do better and more conscientious work for this study.\footnote{“The Explanation of Luther’s Catechism,” by Pontoppidan, which is so highly prized by Norwegian Lutherans, is based on and largely drawn out of this “Explanation” of Spener.} The work of Spener and the Pietists did much to make good catechising more general and more effective. In proportion as their method and spirit have been conserved and followed has the work of catechising been fraught with blessing.

\randnotiz{Position and failure of Rationalism.} The rationalistic age, which cast its withering blight over the whole religious life of Germany, clouding the former clear conception of truth, and freezing out the warmth of personal fellowship and communion with the blessed Saviour, wrought ruin also with the religious instruction of the children of the Church. Rationalism denies the old Bible doctrine of sin. There are no native seeds of evil in man. The seeds and germs of everything that is good are there. The germs of all necessary truth are there. They need not be put in. They need not be revealed. They need only to be drawn out, developed and matured. Hence instruction in religion should not concern itself so much with imparting truth, should not follow the old fogy and
\quellseite{120}
unscientific method of memorizing, explaining and applying what Luther taught, but should, after the Socratic fashion, draw all needed truth out of the child. A truth that cannot be found in and questioned out of a normal child is, if not hurtful, at best, useless. The wiseacres, wise above what is written, forgot two things in their feverish advocacy of the Socratic method. First, that Socrates was dealing with the natural truth of the natural world as it is related to the natural man. But the Christian teacher is dealing with the supernatural truth of the spiritual realm for the spiritual life of man. Second, that Socrates was dealing with mature minds while the Christian Catechist deals with immature children. And, therefore, in spite of the advocacy of Basedow, and Barth, and Dolz, and Dinter, and Herder, and Froebel, and Mosheim, and, to a certain extent, Pestallozzi, the method failed and the fruitage was inexpressibly sad and unsatisfying. The better part of the Church and people were glad to come back to Luther’s Catechism and Spener’s method. \randnotiz{Return to Luther’s Catechism.} And so Luther’s Catechism was welcomed back into the schools and churches of Hesse, Würtemberg and all the Lutheran parts of Germany. In Austria, also, since the edict of toleration of Joseph II., Luther’s Catechism has enjoyed a large and increasing use. In the German parts of Russia and the Baltic provinces, in Protestant
\quellseite{121}
Poland, in the Scandinavian countries and in Schleswig-Holstein it is to-day the principal book for the Christian training of the youth. The “churches of the Augsburg Confession” in France use it. The Moravians use Luther’s Catechism more than any other, following, as we have seen, the example of their illustrious Count Zinzendorf. Before we leave the story of the use of Luther’s Catechism, we must note a custom that is peculiar to Germany and the Germans in connection with catechising. \randnotiz{Kinderlehre.} We refer to the institution called Kinderlehre. Kinderlehre or Christenlehre. This means the catechisation of those already confirmed, by the pastor at a special Sunday service, generally in the afternoon. It generally consists in a free, informal and simple address on a part of the Catechism, usually on one of Luther’s questions. After such a clear and straightforward explanation and application the pastor questions those present as to their understanding of the matter before them. In many places the whole congregation, as well as those confirmed in recent years, are invited. The efficiency, profit and blessing of the service depend upon the Catechist. There are prosy, dull and dry leaders of the Kinderlehre, as well as preachers in the pulpit. And there are those who love the service, put interest and life into it, and make it both attractive and blessed to all. The practice dates back to the
\quellseite{122}
Reformation age, fell into disuse in the age of dry orthodoxism, was revived by Spener, and degenerated again under the influence of Rationalism. Muhlenberg introduced it in the churches which he organized and served.\footnote{See Mann’s “Life and Times of Muhlenberg,” p. 289.} The Pennsylvania Ministerium frequently recommended it, and it is still practiced in many German churches, where it is intended to answer in place of a Sunday school.

Some kind of an adaptation in the form of stated services for the young, with clear, interesting and well illustrated addresses on the truths of the Catechism, encouraging free questions and discussions might prove a blessing in many a small town and village.

\randnotiz{Use by early Lutherans in America.} The history of the use of our Catechism in the English Lutheran churches of America is full of interest, warning and encouragement. As we have seen, the early Lutherans, both those from Halle and those from Sweden, set great store by our Catechism.

They not only prepared editions and explanations suited to their times and surroundings, but were diligent in the regular work of Kinderlehre. As early as 1693 we find a Swedish layman, Charles Springer, sending to Sweden for two hundred copies of the Catechism. In 1736, Pastor Boltzius gives an account of his method of teaching the Catechism: “It was the subject of explanation
\quellseite{123}
at the evening services which he held every day with his people. He inculcated it for the first time among the young, and refreshed the minds of adults. He speaks of the great advantage of this daily instruction, over the weekly, where one part is forgotten before what immediately follows is set forth, and notes the deep interest of the people, who give the same attendance on these exercises as is accorded the regular Sunday services. He was always careful to limit the entire service to one-half hour.

“‘Last evening,’ he writes, ‘we reviewed the Seventh and catechised upon the Eighth Commandment. The consciences of some in the congregation were deeply moved. In the exposition of the Commandments to simple hearers we find it highly necessary to indicate the sins forbidden therein, and the virtues enjoined not only generally, but they must be clearly specified according to the circumstances of the hearers.’

“Another entry runs: ‘One of those who was recently comforted by his confession of several sins against the Seventh Commandment, brought to me some money to-day, to restore a part of his unrighteous gain.’”\footnote{Jacobs’ “History of Lutheran Church in United States,” p. 165.}

In later generations our Church, which had made the fatal mistake of not planting her own institutions
\quellseite{124}
of learning, became affected by the rationalism and the emotionalism that abounded in the American Reformed churches. Luther’s Catechism fell into disfavor and disuse. \randnotiz{Era of disfavor and disuse} Substitutes were offered in its place. So-called explanations, which explained away Luther’s sense, were circulated. Such a substitute for Luther’s was Quitman’s Catechism, published “with consent and approbation” of the New York Ministerium in 1814. It was rationalistic throughout, and removed the very foundations as well as the essence of the faith of the Church of the Reformation. In circles where there was spirituality enough to be repelled by such subtle and soul-destroying forms of unbelief, emotionalism under the form of so-called “New Measurism” or Revivalism, took the place of catechisation. Indifference to doctrine fostered indifference if not hostility to catechisation. In many places in the Lutheran churches, especially those liberal English Lutheran churches that wanted to be like the other churches around them, it went entirely out of fashion. It was no rare thing to find people who had grown up in an English Lutheran church who had not seen a copy of Luther’s Catechism. The results, as may be imagined, were disastrous. When people know of no difference between the Lutheran Church and others, why should they love or sacrifice for her? \randnotiz{Results} If one is as good as another it is foolish,
\quellseite{125}
it is wasteful, if not sinful, to build up Lutheran missions where there are other churches.

It took some liberal Lutherans a long time to learn these self-evident truths. Some, indeed, have not learned them to this day. This diminishing class is careless as to thorough catechising. They do not build up strong Lutheran churches. If all would accept their principles and follow their methods, the Lutheran Church would have no future. And she would deserve none.

But we are glad to know that these liberals are not the Lutheran Church. Wherever she is true to herself, she insists that her pastors be good Catechists. \randnotiz{Demand and desire of the true Lutheran Church} She desires that her children and youth shall be thoroughly and heartily instructed in Luther’s Small Catechism, she knows that her future is bound up with the proper instruction and training of her youth. For this she knows no better text-book than Luther’s Small Catechism. As our Church is faithful or unfaithful to the teachings of this book will she be strong or weak. When she no longer teaches or learns its lessons her history will be at an end. She will be dead.
